% paper4-main.tex — The Topology Programmer
% TASUMER MAF Paper 4: Cybernetic feedback as general-purpose geometric prior
%
% ============================================================================
% ⚠ AUDIT BLOCKERS — DO NOT SUBMIT / COMPILE-FOR-RELEASE UNTIL RESOLVED
% (added 2026-07-04 from the controller-exposure classification;
%  full worklist: docs/PAPER4_EXPOSURE_CLASSIFICATION.md §9)
%   1. RETRACTED DATA: line ~218 cites "Wan 2.1 14B" as a model scale — that is
%      the Holly Battery run, retracted 2026-03-13 (dataset provenance / L-026).
%      The 14B scale claim must be dropped or replaced with a clean run.
%   2. FALSE ATTRIBUTION: section8_discussion.tex ~L214-216 attributes the 36.2%
%      P0 overhead to "control law evaluation"; the P0 run (train_comparison.py)
%      has NO Steersman. The 36.2% number is clean; the causal sentence is false.
%   3. NUMBER/CONSISTENCY: tab:attractor H-ch4 = 0.0836 is a step-1100 PARTIAL,
%      not the 0.0918 final; "34 matching checkpoints" (L185) vs "6 matching
%      steps" (§4.3) for r=1.0000; spectral-attractor Fiedler ~0.09 is at BM-000
%      null percentile 17 (near-init noise) — reframe, do not cite as structure.
%   Numbers otherwise frozen pending the C6b controller-stability fix.
% ============================================================================
%
% Compile: pdflatex paper4-main && bibtex paper4-main && pdflatex paper4-main && pdflatex paper4-main

\documentclass[11pt]{article}

% --- packages ---
\usepackage[margin=1in]{geometry}
\usepackage{amsmath,amssymb,amsfonts}
\usepackage{booktabs}
\usepackage{graphicx}
\usepackage{hyperref}
\usepackage{xcolor}
\usepackage{multirow}
\usepackage{float}
\usepackage[numbers]{natbib}
\usepackage{algorithm}
\usepackage{algorithmic}
% \usepackage{microtype}  % disabled: MiKTeX font expansion issue
\usepackage{bookmark}

\hypersetup{
  colorlinks=true,
  linkcolor=blue!60!black,
  citecolor=blue!60!black,
  urlcolor=blue!60!black,
  hypertexnames=false
}

% --- metadata ---
\title{%
  The Topology Programmer:\\
  Cybernetic Feedback as a General-Purpose Geometric\\
  Prior for Neural Network Adapters%
}

\author{%
  Timothy Paul Bielec%
  \thanks{Corresponding author. \texttt{timothy@promptcrafted.com}}%
  \\
  Promptcrafted LLC / TASUMER MAF\\
  Los Angeles, CA
}

\date{March 2026}

\begin{document}
\maketitle

% ===================================================================
\begin{abstract}
A companion paper established that cybernetic feedback drives
block-diagonal emergence in 6-channel LoRA bridges aligned with
rhombic dodecahedral (RD) face-pair geometry. This paper asks whether
the mechanism is specific to the RD or general-purpose, and
characterizes its boundaries. The answer: the Steersman is a
\emph{selective geometric compiler}. It programs any pair specification
derived from polytope co-axial structure and rejects all others. The
same feedback mechanism, with no modification beyond the pair
specification, programs 4-dimensional tesseract geometry into 8-channel
bridges (four independent $2 \times 2$ blocks, co-axial/cross-axial
coupling ratio \textbf{41{,}564:1}, Fiedler eigenvalue 0.000191, clean
$4{+}4$ eigenvalue split) and octahedral geometry into 4-channel bridges
(two $2 \times 2$ blocks, \textbf{473{,}622:1}---the programme's
strongest signal). Four negative controls sharpen the claim:
spectral-only training produces a universal attractor at Fiedler
${\approx}\,0.09$ across $n \in \{3, 4, 8, 12\}$; wrong-labels
(random pair partition) and resonance (number-theoretic pairs) both
collapse to co/cross ${\sim}\,10^{-5}$:1, functionally
indistinguishable from each other ($<10\%$ divergence across 10{,}000
steps); and an emanation architecture (hierarchical master bridge)
produces weak coherence (1.12:1) without block-diagonal structure.
These outcomes define four cleanly separated dynamical regimes:
block-diagonal ($\rho > 40{,}000$:1), spectral attractor
($\text{Fiedler} \approx 0.09$), hierarchical coherence
($\rho \approx 1$:1), and connectivity collapse
($\rho \approx 10^{-5}$:1). The requirement is spatial coherence:
number-theoretic structure without geometric embedding is functionally
equivalent to random noise. Tesseract training is reproducible
($r = 1.0000$ across independent runs). The bridge is a programmable
topology substrate; the pair specification is the program; the
Steersman compiles valid geometric programs and rejects incoherent ones.
\end{abstract}

\noindent\textbf{Keywords:} topology programming, cybernetic feedback,
bridge matrix, selective geometric compiler, spectral attractor,
multi-channel LoRA, parameter-efficient fine-tuning, block-diagonal
structure


% ===================================================================
\section{Introduction}
\label{sec:intro}

Low-rank adaptation~\cite{hu2022lora} modifies pretrained models by
injecting rank-$r$ update matrices $\Delta W = BA$ alongside frozen
weight matrices. Multi-channel extensions~\cite{bielec2026weights,
bielec2026bridge} partition the rank dimension into $n$ channels with
an $n \times n$ bridge matrix governing inter-channel coupling. The
bridge adds $n^2$ learnable parameters (36 at $n = 6$) and introduces
a new degree of freedom: the \emph{topology} of the adapter.

Paper~3~\cite{bielec2026bridge} demonstrated that cybernetic
feedback---a closed-loop control system called the Steersman---drives
the bridge toward 3-axis block-diagonal structure aligned with the
rhombic dodecahedron's face-pair geometry. The result was consistent:
100\% block-diagonal across six experiments, 42{,}500+ bridges, two
model scales, and three initialization strategies, with 0\% in all
non-cybernetic controls. Co-planar/cross-planar coupling ratios
exceeded 22{,}000:1. The Bridge Fiedler eigenvalue converged to
${\sim}10^{-5}$, indicating near-perfect disconnection between blocks.

That paper left a question open. The Steersman encodes
rhombic dodecahedral geometry through a specific pair specification:
six channels decompose into three co-planar pairs aligned with the
$x$, $y$, and $z$ axes. Block-diagonal emergence with three
independent $2 \times 2$ blocks is the unique structural solution for
this specification. Is this result specific to the RD---a property of
the 3-axis geometry, the six-channel count, or the particular
face-pair decomposition---or is the Steersman a general-purpose
mechanism for programming arbitrary topology into learnable bridge
matrices?

This paper provides the answer: the Steersman is general, but
\emph{selective}. With no modification to the feedback mechanism,
control laws, or training procedure---only the pair specification
changes---the Steersman programs 4-dimensional tesseract geometry
into 8-channel bridges (T-001r2: co-axial/cross-axial coupling ratio
\textbf{41{,}564:1}, Fiedler 0.000191, clean $4{+}4$ eigenvalue
split) and octahedral geometry into 4-channel bridges (O-001:
per-bridge mean \textbf{473{,}622:1}---the programme's strongest
signal).

Four negative controls define the boundaries. Spectral-only training
across $n \in \{3, 4, 8, 12\}$ reveals a universal spectral attractor
at Fiedler ${\approx}\,0.09$---the bridge's natural connectivity
equilibrium when the Steersman provides spectral pressure but no
topological direction. Wrong-labels (WL-001, random pair partition)
and resonance (R-001, number-theoretic pairs) both collapse to
co/cross ${\sim}\,10^{-5}$:1 with functionally identical trajectories
($<10\%$ divergence), demonstrating that algebraic structure without
spatial embedding carries no information the contrastive loss can use.
An emanation architecture (E-001, hierarchical master bridge with
per-layer offsets) produces weak directional coherence (1.12:1)
without block-diagonal structure, establishing that topology
programming requires direct contrastive feedback to individual bridge
matrices.

These outcomes define four cleanly separated dynamical regimes
with no intermediate or overlapping cases: (1)~block-diagonal
($\rho > 40{,}000$:1), (2)~spectral attractor ($\text{Fiedler}
\approx 0.09$, $\rho \approx 1$:1), (3)~hierarchical coherence
($\rho \approx 1$:1, Fiedler $\approx 0.08$), and (4)~connectivity
collapse ($\rho \approx 10^{-5}$:1, Fiedler $\approx 10^{-5}$).

We make the following contributions:

\begin{enumerate}
  \item \textbf{4D topology programming.} The Steersman programs
    tesseract geometry ($n = 8$, four co-axial pairs) with the same
    fidelity as 3D rhombic dodecahedral geometry ($n = 6$, three
    co-planar pairs). Co-axial/cross-axial ratio: 41{,}564:1 at 10K
    steps. Fiedler: 0.000191. Clean $4{+}4$ eigenvalue split.
  \item \textbf{Minimum-geometry programming.} Octahedral contrastive
    ($n = 4$, two co-axial pairs) produces the programme's strongest
    signal: 473{,}622:1 per-bridge mean coupling ratio.
  \item \textbf{Spectral attractor universality.} Spectral-only
    training converges to Fiedler ${\approx}\,0.09$--$0.10$ across
    $n \in \{3, 4, 8, 12\}$---a 10.4\% band across a $4\times$ range
    of channel counts.
  \item \textbf{The selective geometric compiler.} The Steersman
    compiles polytope-derived pair specifications into block-diagonal
    structure and rejects non-spatial specifications (random partitions,
    number-theoretic relationships) by driving the bridge into
    connectivity collapse.
  \item \textbf{Number-theoretic equivalence to random noise.}
    Wrong-labels (random partition) and resonance (number-theoretic
    pairs) produce functionally identical collapse trajectories
    ($<10\%$ divergence across 10{,}000 steps).
  \item \textbf{Four-regime taxonomy.} The complete experimental
    programme produces four cleanly separated dynamical regimes
    with no intermediate outcomes.
  \item \textbf{Reproducibility.} $n = 8$ tesseract contrastive
    replicates with Pearson $r = 1.0000$ across 34 matching checkpoints.
\end{enumerate}

The remainder of this paper is organized as follows.
Section~\ref{sec:background} reviews Paper~3's findings and related
work on programmable inductive biases.
Section~\ref{sec:method} introduces the generalized pair specification
interface and the five experimental topologies.
Section~\ref{sec:tesseract} presents the 4D programming results.
Section~\ref{sec:attractor} characterizes the spectral attractor.
Section~\ref{sec:programmability} tests topology programmability
through wrong-labels and octahedral controls.
Section~\ref{sec:nongeometric} examines non-geometric topologies
(resonance and emanation) and presents the four-regime taxonomy.
Section~\ref{sec:discussion} discusses implications and limitations.
Section~\ref{sec:conclusion} concludes.


% ===================================================================
\section{Background}
\label{sec:background}

\subsection{Paper 3: The Learnable Bridge}

The companion paper~\cite{bielec2026bridge} established three results
relevant to the present work.

First, the Steersman---a cybernetic feedback mechanism applying
contrastive and spectral pressure to multi-channel LoRA bridge
matrices---produces 100\% block-diagonal structure at $n = 6$
channels, with three independent $2 \times 2$ blocks aligned to the
rhombic dodecahedron's three coordinate axes. The co-planar/cross-planar
coupling ratio exceeded 82{,}154:1 across 60{,}000+ bridges and three
model scales (TinyLlama~1.1B, Qwen~7B, Wan~2.1~14B), while all
non-cybernetic controls produced 0\% block-diagonal structure.

Second, the block-diagonal attractor is initialization-independent:
identity, geometric bias, and corpus-coupled (I~Ching complementary
trigram) initializations converge to the same block-diagonal structure
within 200~steps (Figure~\ref{fig:fi002_init_independence}). The
I~Ching initialization, which actively opposes the target topology,
is suppressed 99.5\% in 900~steps.

Third, the channel-count ablation at $n \in \{3, 4, 6, 8\}$ showed
that the contrastive loss encoding RD face-pair geometry is necessary
and sufficient for block-diagonal emergence at $n = 6$; at other
channel counts (where contrastive loss was not defined), spectral
regularization alone produces uniform coupling with Fiedler eigenvalue
${\sim}0.09$. Paper~3 interpreted this convergence as the spectral
loss driving connectivity without topological direction.

The present paper tests this interpretation directly. By defining
contrastive loss for additional geometries ($n = 8$ tesseract, $n = 4$
octahedron) and additional non-geometric specifications ($n = 6$
wrong-labels, resonance), we characterize the full boundary between
successful and unsuccessful topology programming.

\subsection{Programmable Inductive Biases}

The idea that geometric or structural priors improve neural network
training is well-established. Convolutional networks encode
translation equivariance~\cite{lecun1998}. Graph neural networks
encode permutation equivariance over relational
structures~\cite{scarselli2009, kipf2017}. Geometric deep
learning~\cite{bronstein2021} provides a unified framework for
equivariant architectures over groups and manifolds.

These priors are typically \emph{architectural}: the network's
connectivity is designed to enforce the desired invariance. The
Steersman operates differently. It does not constrain the bridge
architecture---the bridge remains a dense, unconstrained
$n \times n$ matrix throughout training. Instead, it applies
\emph{training-time feedback} that guides the bridge toward a target
topology through the optimization dynamics. The structure is
discovered, not imposed.

Hypernetworks~\cite{ha2017} generate network weights conditioned on
external inputs, enabling task-conditional structure. Meta-learned
optimizers~\cite{andrychowicz2016} adapt optimization dynamics
through learned feedback. The Steersman differs from both: its
control laws are fixed (not learned), its state space is
low-dimensional (three scalar diagnostics), and its target topology
is specified declaratively through pair indices rather than learned
from data. This simplicity is a design choice: we aim to demonstrate
that minimal feedback with explicit geometric specification is
sufficient for robust topology programming.

\subsection{Tesseract Geometry}

The tesseract (4-cube, hypercube, $\gamma_4$) is the
four-dimensional analogue of the cube. It has 16~vertices, 32~edges,
24~square faces, and 8~cubic cells. For the present work, the
relevant property is its axis structure: the 16~vertices decompose
into 4~co-axial pairs along the four coordinate axes of
$\mathbb{R}^4$. This decomposition provides the pair specification
for 8-channel bridges: each of the four coordinate axes maps to one
co-axial channel pair, yielding four independent $2 \times 2$ blocks
if the topology is successfully programmed.

Higher-dimensional geometric structure has appeared in ML in several
contexts: hyperbolic embeddings for hierarchical
data~\cite{nickel2017}, algebraic geometry for neural network loss
landscapes~\cite{watanabe2009}, and quaternionic
networks~\cite{parcollet2019} for 4D rotation representation. The
present work uses 4D geometry differently: not as a representation
space or an architectural constraint, but as a target topology for
a feedback-driven bridge matrix.

\subsection{Spectral Graph Theory and Bridge Structure}

The bridge matrix $\mathbf{M} \in \mathbb{R}^{n \times n}$ can be
interpreted as a weighted adjacency matrix of a complete graph on
$n$~nodes. The Fiedler eigenvalue~\cite{fiedler1973}---the
second-smallest eigenvalue of the weighted graph Laplacian---measures
algebraic connectivity: how well-connected the graph is. A Fiedler
value near zero indicates near-disconnection; a positive value
indicates that all nodes participate in a connected component.

This spectral perspective provides the diagnostic framework for the
Steersman's control laws. The contrastive loss operates on the
distinction between co-axial entries (to be strengthened) and
cross-axial entries (to be suppressed). When this distinction
produces block-diagonal structure, the bridge graph decomposes into
disconnected cliques (the blocks), and the Fiedler eigenvalue
approaches zero. When the distinction is absent (spectral-only) or
incoherent (wrong-labels, resonance), the Fiedler eigenvalue reflects
the outcome: stable connectivity (${\sim}0.09$) or total collapse
(${\sim}10^{-5}$).

The stochastic block model~\cite{holland1983} provides the graph-theoretic
analogue of the Steersman's output: a random graph with known
community structure, where within-community edge probabilities exceed
between-community probabilities. Block-diagonal bridge structure is
the extreme case: within-block coupling is maximal and between-block
coupling approaches zero.


% ===================================================================
% Sections 3--8 are modular files
% ===================================================================
\section{Method: Generalized Pair Specification}
\label{sec:method}

The Steersman's contrastive loss, as defined in
Paper~3~\cite{bielec2026bridge}, encodes the rhombic dodecahedron's
face-pair geometry through a fixed partition of channel indices into
co-planar and cross-planar sets. We generalize this mechanism by
replacing the hard-coded RD partition with a configurable pair
specification interface. Five experimental topologies test the
boundaries of this generalization.

\subsection{The Pair Specification Interface}
\label{sec:pairs}

The pair specification is a function
$\texttt{pairs}(n) \to \{(i_1, j_1), \ldots, (i_k, j_k)\}$ that
returns a set of \emph{co-axial} channel index pairs for a given
channel count~$n$. All remaining pairs of distinct indices are
classified as \emph{cross-axial}. The contrastive loss operates on
this partition identically to Paper~3: it rewards co-axial coupling
and penalizes cross-axial coupling. The Steersman's three control
laws (connectivity trending, directionality stagnation, stability
monitoring) operate without modification.

For $n = 6$, the RD specification is:
\begin{equation}
  \texttt{pairs}_{\text{RD}}(6) = \{(0,5),\; (1,4),\; (2,3)\}
  \label{eq:rd_pairs}
\end{equation}
These are the three face-pair axes of the rhombic dodecahedron. Each
pair shares a vanishing coordinate in the face-normal vectors,
producing three independent $2 \times 2$ blocks in the bridge.

For $n = 8$, the tesseract specification is:
\begin{equation}
  \texttt{pairs}_{\text{tess}}(8) = \{(0,1),\; (2,3),\; (4,5),\;
    (6,7)\}
  \label{eq:tess_pairs}
\end{equation}
These are the four co-axial vertex pairs of the tesseract, aligned
with the four coordinate axes of $\mathbb{R}^4$.

For $n = 4$, the octahedral specification is:
\begin{equation}
  \texttt{pairs}_{\text{oct}}(4) = \{(0,1),\; (2,3)\}
  \label{eq:oct_pairs}
\end{equation}
These are two of the octahedron's three vertex-opposition axes,
selected to match the channel count. The octahedron is the simplest
regular polytope admitting a pair decomposition and provides the
minimum-channel-count test of the geometric coherence hypothesis.

The interface imposes no constraint on which indices are paired.
Any partition of $n$~channels into $\lfloor n/2 \rfloor$ pairs
constitutes a valid specification. This generality enables the
non-geometric controls described below.

\subsection{Tesseract Topology ($n = 8$)}
\label{sec:tesseract_method}

The tesseract (4-cube, $\gamma_4$) has 16~vertices decomposing
into 4~co-axial pairs along the four coordinate axes of
$\mathbb{R}^4$. The pair specification
$\texttt{pairs}_{\text{tess}}(8)$ assigns one channel pair per
tesseract axis. If the Steersman generalizes beyond 3D, this
specification should produce four independent $2 \times 2$
blocks---a $4{+}4$ eigenvalue split.

\subsection{Octahedral Topology ($n = 4$)}
\label{sec:octahedral_method}

The octahedron has 6~vertices forming 3~co-axial pairs. At $n = 4$
channels, two of these axes are specified:
$\texttt{pairs}_{\text{oct}}(4) = \{(0,3),\; (1,2)\}$.
This represents the minimum geometry: only 2~co-axial pairs and
4~cross-axial pairs. The cross-axial/co-axial pair count ratio
($4/2 = 2$) is smaller than for the tesseract ($24/4 = 6$) or the
RD ($12/3 = 4$), providing the strongest co-axial signal per
channel pair.

\subsection{Wrong-Labels Topology ($n = 6$)}
\label{sec:wrong_labels_method}

To test whether the Steersman programs whatever topology it is told
or whether geometrically correct pair specifications are required,
we define a \emph{wrong-labels} specification: $n = 6$ channels
partitioned into three pairs by random assignment rather than RD
face-pair geometry. The Steersman's contrastive loss treats these
random pairs as co-axial and all other pairs as cross-axial. If the
Steersman is a general topology programmer accepting arbitrary
partitions, it should produce block-diagonal structure aligned with
the random pairs. If geometric coherence is required, the random
partition should fail.

\subsection{Circular Resonance Topology ($n = 6$)}
\label{sec:resonance_method}

The resonance topology derives pair specifications from
number-theoretic relationships among the tracked prime threads of the
companion research programme~\cite{bielec2026shape}. For $n = 6$, the
pairs are assigned by algebraic relationships: Sophie Germain pairing
($2(11) + 1 = 23 \to$ channels $(0, 2)$), consecutive-prime pairing
($29, 31 \to$ channels $(3, 4)$), and residue pairing (same residue
$\bmod\,6 \to$ channels $(1, 5)$). This specification tests whether
algebraic structure---without spatial embedding---can serve as a
topology specification.

\subsection{Emanation Architecture ($n = 6$)}
\label{sec:emanation_method}

The emanation architecture introduces hierarchical bridge structure.
Rather than independent bridge matrices per adapter (as in all prior
experiments), the emanation design maintains a single \emph{master
bridge} shared across all layers, with per-layer \emph{offset
matrices} that capture local variation:
\begin{equation}
  \mathbf{M}^{(l)} = \mathbf{M}_{\text{master}} +
    \boldsymbol{\Delta}^{(l)}
  \label{eq:emanation}
\end{equation}
where $l$ indexes the transformer layer and offsets are scaled by
$1/\sqrt{l}$. The Steersman monitors a \emph{coherence} metric---the
mean Frobenius norm of the offsets relative to the master---and
applies regularization pressure when individual layers fragment
excessively. The same RD-derived co-axial pairs are specified, but
the contrastive loss operates through the master matrix rather than
per-layer bridges directly.

\subsection{Training Setup}

All experiments use TinyLlama-1.1B-Chat with the same hyperparameters
as Paper~3: learning rate $2 \times 10^{-4}$, batch size~2, gradient
accumulation~8 (effective batch size~16), LoRA rank~24, identity
initialization, random seed~42, maximum sequence length~512,
alpaca-cleaned dataset, 10{,}000 training steps. The channel size
$s = r/n$ varies with~$n$: $s = 8$ ($n = 3$), $s = 6$ ($n = 4$),
$s = 4$ ($n = 6$), $s = 3$ ($n = 8$).

Table~\ref{tab:topology_specs} summarizes all pair specifications.

\begin{table}[H]
\centering
\caption{Pair specifications for all experimental topologies. Co-axial
  pairs receive positive contrastive gradient; all remaining
  off-diagonal index pairs are cross-axial and receive negative
  gradient. $|\mathcal{C}|$ = co-axial pair count,
  $|\mathcal{X}|$ = cross-axial pair count.}
\label{tab:topology_specs}
\small
\begin{tabular}{llclcc}
\toprule
Topology & $n$ & Co-axial pairs & Source & $|\mathcal{C}|$ & $|\mathcal{X}|$ \\
\midrule
RD (3D) & 6 & $(0,5), (1,4), (2,3)$ & Face opposition & 3 & 12 \\
Tesseract (4D) & 8 & $(0,1), (2,3), (4,5), (6,7)$ & Vertex opposition & 4 & 24 \\
Octahedral (3D) & 4 & $(0,1), (2,3)$ & Vertex opposition & 2 & 4 \\
Wrong-labels & 6 & Random partition & None (control) & 3 & 12 \\
Resonance & 6 & $(0,2), (3,4), (1,5)$ & Number-theoretic & 3 & 12 \\
\bottomrule
\end{tabular}
\end{table}

% ===================================================================
\section{4D Topology Programming}
\label{sec:tesseract}

\subsection{Experimental Design}

The tesseract experiment (T-001) uses $n = 8$ channels with pair
specification $\texttt{pairs}_{\text{tess}}(8)$
(Eq.~\ref{eq:tess_pairs}). Two predictions frame the result:

\begin{itemize}
  \item \textbf{Prediction A (4D programming works):} The bridge
    develops four independent $2 \times 2$ blocks, with co-axial
    coupling concentrated within each tesseract axis pair and
    cross-axial coupling suppressed. The eigenvalue spectrum shows a
    $4{+}4$ split: four eigenvalues near zero (suppressed cross-axial
    connections) and four elevated (active co-axial connections).
  \item \textbf{Prediction B (3D attractor):} The bridge collapses to
    three dominant blocks regardless of the 4-pair specification.
    Three-dimensional geometry would be an intrinsic structural limit,
    and the Steersman would be RD-specific rather than general.
\end{itemize}

Training uses TinyLlama-1.1B with identity initialization and
channel size $s = 3$ ($24 / 8$). Four same-seed runs were conducted
(Section~\ref{sec:reproducibility}): T-001r1 (original; crashed at
${\sim}$step~2{,}700 during a March-2026 machine failure, with its
first 1{,}600~steps recovered from its training log), T-001r1b (a
fresh same-seed relaunch, 7{,}100~steps), T-001r2 (complete,
10{,}000~steps), and T-001r3 (a fourth complete run, 10{,}000~steps,
executed on separate hardware).
% Dated edit 2026-07-10: T-001r3 (fourth run) added as the fourth same-seed
% replicate. Source: results/T-001-full-r3/results.json + EXPERIMENT_TRACKER
% 2026-07 sprint update (L-006 dated addition).

\subsection{Results: Prediction A Confirmed}

The Steersman programs 4D tesseract geometry.
Table~\ref{tab:tesseract} reports the primary metrics.

\begin{table}[H]
\centering
\caption{Tesseract contrastive results ($n = 8$). $\rho$ = co-axial
  to cross-axial coupling ratio. Fiedler = Bridge Fiedler eigenvalue
  (mean across adapters). Val loss reported at final step.}
\label{tab:tesseract}
\small
\begin{tabular}{lcccccc}
\toprule
Run & Steps & Fiedler & $\rho$ & Val Loss & Eigenvalue Split & BD \\
\midrule
T-001r1$^{\dagger}$ & 1{,}600 & 0.00055 & 4{,}729:1 & --- & --- & --- \\
T-001r1b & 7{,}100 & 0.000362 & 20{,}944:1 & 0.4049 & $4{+}4$ & Yes \\
T-001r2 & \textbf{10{,}000} & \textbf{0.000191} & \textbf{41{,}564:1} & \textbf{0.4016} & $4{+}4$ & \textbf{Yes} \\
% Dated edit 2026-07-10: T-001r3 fourth run (Hermes RTX 4090); endpoint from
% results/T-001-full-r3/results.json checkpoints[-1] (step 10000).
T-001r3 & 10{,}000 & 0.000193 & 41{,}812:1 & 0.4017 & $4{+}4$ & Yes \\
\bottomrule
\end{tabular}

\smallskip
\raggedright\footnotesize $^{\dagger}$~T-001r1's results file was
overwritten by the T-001r1b relaunch (Section~\ref{sec:reproducibility});
values shown are the last recovered from its training log (step~1{,}600).
Its contemporaneously recorded endpoint (5{,}395:1 at step~2{,}700) has
no surviving artifact and is not tabulated.
\end{table}

Both completed runs produce four independent $2 \times 2$ blocks. The
eigenvalue spectrum exhibits a clean $4{+}4$ split: four eigenvalues
suppressed to $\mathcal{O}(10^{-5})$ (the cross-axial connections)
and four elevated to ${\sim}1.5$--$1.6$ (the co-axial connections within
each block). This is the structural signature of 4D block-diagonal
geometry: each tesseract axis operates as an independent channel
pair, coupled internally and decoupled from the other three axes.

The coupling ratio trajectory rises by three orders of magnitude:
1{,}931:1 at step~400, 5{,}079:1 at step~2{,}800, and \textbf{41{,}564:1}
at step~10{,}000 (T-001r2). Bridge Fiedler at 0.000191 confirms the topological
bifurcation: the spectral attractor at ${\sim}0.09$
(Section~\ref{sec:attractor}) has been broken by the 4D contrastive
specification, collapsing Fiedler by approximately $494\times$
(from $0.0944$ at $n = 8$ spectral-only to $0.000191$). This
collapse is of the same order as the $1{,}020\times$ bifurcation
at $n = 6$ ($0.0918 \to 0.00009$).

Validation loss at 0.4016 is within 0.17\% of the baseline,
consistent with Paper~3's finding that topology programming does not
degrade task performance. The lm-eval benchmark suite confirms this
at the downstream level: the original T-001r1 adapter (2{,}700~steps,
merged and evaluated on 2026-03-16, before the relaunch overwrote its
source artifact; the evaluation outputs survive) shows a mean delta of
$-0.75\%$ across six 0-shot tasks
(HellaSwag, ARC-Easy, PIQA, WinoGrande, OpenBookQA, BoolQ), with
mixed-direction changes (three tasks slightly down, two slightly up,
one negligible). Bridge topology is orthogonal to task capability.

Prediction~A is confirmed: 4D programming works. The Steersman
programs the dimensionality it is told to program.

\subsection{Reproducibility: Four Same-Seed Runs}
\label{sec:reproducibility}

% Dated edit 2026-07-10: run count three -> four; T-001r3 endpoint paragraph
% added below. Sources: results/T-001-full-r3/results.json checkpoints[-1];
% rel dev vs r2's 41,564 per results/EXPERIMENT_TRACKER.md 2026-07 sprint update.
Four runs of the $n = 8$ tesseract contrastive experiment were
conducted with identical hyperparameters and random seed. The
provenance is disclosed in full. T-001r1, the original run, crashed
at ${\sim}$step~2{,}700 during a March-2026 machine failure; when the
experiment was relaunched fresh under the same seed (T-001r1b,
reaching 7{,}100~steps), the relaunch overwrote T-001r1's results
file. T-001r1's first 1{,}600~steps were subsequently recovered from
its surviving training log. T-001r2 ran to completion
(10{,}000~steps). A fourth run, T-001r3, was executed later on
separate hardware (Appendix~\ref{app:details}), also running to
completion (10{,}000~steps) and written to a fresh output directory so
that no existing artifact was overwritten. The first three
trajectories, which carry recoverable per-step feedback logs,
constitute independent same-seed replicates (one partially recovered)
and are cross-correlated below; T-001r3 provides an independent
endpoint confirmation.

\begin{table}[H]
\centering
\caption{Reproducibility of tesseract contrastive training: pairwise
  trajectory agreement between the three same-seed runs with
  recoverable per-step feedback logs, at matched
  feedback steps. T-001r1 pairs span its recovered range
  (steps 200--1{,}600); the T-001r1b~$\times$~T-001r2 pair spans
  their full overlap.}
\label{tab:reproducibility}
\small
\begin{tabular}{lcccc}
\toprule
Pair & $\rho$ Pearson $r$ & $\rho$ max dev & Fiedler $r$ & Steps \\
\midrule
T-001r1 $\times$ T-001r1b & 0.9952 & 10.4\% & 0.9994 & 15 \\
T-001r1 $\times$ T-001r2  & 0.9975 & 5.2\%  & 0.9996 & 15 \\
T-001r1b $\times$ T-001r2 & 0.9984 & 7.2\%  & 0.9999 & 71 \\
\bottomrule
\end{tabular}
\end{table}

Table~\ref{tab:reproducibility} reports the pairwise agreement. All
three pairings track at co-axial/cross-axial ratio correlation
$r \geq 0.995$ and Fiedler correlation $r \geq 0.9994$, with maximum
relative deviations of up to ${\sim}10\%$ at matched steps. The
deviations are the expected signature of floating-point
non-associativity amplified over thousands of optimizer steps; the
correlations show that the topology-programming
\emph{dynamics}---the trajectory shape, the three-order-of-magnitude
ratio growth, the Fiedler collapse---are reproducible across
independent executions. Notably, the run pair that never coexisted
(the recovered T-001r1 and its replacement T-001r1b) agrees nearly as
closely as the pairs that did, confirming that the recovered log
fragment is a faithful record of a genuine same-seed replicate.

T-001r2 continues beyond both earlier runs' termination points,
reaching $\rho = 41{,}564$:1 at step~10{,}000---double T-001r1b's
20{,}944:1 at step~7{,}100---demonstrating continued growth in trend
through the full training run.

% Dated edit 2026-07-10: fourth same-seed run T-001r3 endpoint confirmation.
% Endpoint from results/T-001-full-r3/results.json checkpoints[-1] (co/cross
% 41811.71, Fiedler 1.9347e-4, val 0.4017); rel dev vs r2's 41,564 = 0.596%
% per results/EXPERIMENT_TRACKER.md 2026-07 sprint update.
A fourth same-seed run, T-001r3, was completed on 2026-07-08 on
separate hardware---the Hermes RTX~4090, whereas T-001r2 ran on the
local RTX~6000 Ada (Appendix~\ref{app:details}). Run to the full
10{,}000~steps in a fresh output directory, it reaches
$\rho = 41{,}812$:1 with Bridge Fiedler $1.93 \times 10^{-4}$ at the
final step, reproducing T-001r2's step-10{,}000 endpoint
(41{,}564:1) to within $0.596\%$ relative deviation. That two
independent same-seed executions on \emph{different} GPUs converge on
the coupling ratio to better than one part in 150 is a stronger
reproducibility statement than same-hardware agreement alone: the
programmed 4D block-diagonal endpoint is a property of the pair
specification and the optimizer, not of a particular device's
floating-point environment.

\subsection{Comparison: RD vs.\ Tesseract}

Table~\ref{tab:rd_vs_tess} compares the two confirmed polytope
topologies. Both produce extreme coupling ratios with sub-attractor
Fiedler values and clean eigenvalue splits matching the specified
dimensionality.

\begin{table}[H]
\centering
\caption{Comparison of confirmed polytope topologies. Both use
  contrastive loss with polytope-derived co-axial pair
  specifications. The RD result (H-ch6) is from Paper~3.}
\label{tab:rd_vs_tess}
\small
\begin{tabular}{lcccccc}
\toprule
Geometry & $n$ & Pairs & $\rho$ & Fiedler & Split & Val Loss \\
\midrule
RD (3D) & 6 & 3 co / 12 cross & 70{,}404:1 & 0.00009 & $3{+}3$ & 0.4015 \\
Tesseract (4D) & 8 & 4 co / 24 cross & 41{,}564:1 & 0.000191 & $4{+}4$ & 0.4016 \\
\bottomrule
\end{tabular}
\end{table}

The RD achieves higher $\rho$ (70{,}404 vs.\ 41{,}564) and lower
Fiedler ($0.00009$ vs.\ $0.000191$), likely reflecting the more
favorable co-axial/cross-axial pair ratio at $n = 6$ ($3/12 = 0.25$)
vs.\ $n = 8$ ($4/24 = 0.167$): fewer cross-axial entries to suppress
at $n = 6$ means faster convergence to the block-diagonal attractor.
Both achieve validation loss within 0.03\% of each other and within
0.17\% of the spectral-only baseline.

% ===================================================================
\section{The Spectral Attractor}
\label{sec:attractor}

\subsection{Universal Convergence}

Paper~3 reported that spectral-only training (contrastive loss
disabled) at $n \in \{3, 4, 8\}$ converges to Bridge Fiedler
${\sim}0.09$. We extend the ablation to $n = 12$ and report the
complete results in Table~\ref{tab:attractor}.

\begin{table}[H]
\centering
\caption{Spectral attractor: Bridge Fiedler eigenvalue under
  spectral-only training (no contrastive loss) across channel
  counts. All experiments use TinyLlama-1.1B, identity
  initialization, and identical hyperparameters. H-ch6 (H3) is
  the $n = 6$ experiment with contrastive loss active, shown for
  bifurcation comparison.}
\label{tab:attractor}
\small
\begin{tabular}{lcccccl}
\toprule
ID & $n$ & Contrastive & Steps & Fiedler & $\rho$ & BD \\
\midrule
H-ch3  & 3  & No  & 10{,}000 & 0.0951  & N/A       & No \\
H-ch4  & 4  & No  & 10{,}000 & 0.0836  & ${\sim}1$:1 & No \\
H-ch8  & 8  & No  & 10{,}000 & 0.0944  & ${\sim}1$:1 & No \\
H-ch12 & 12 & No  & 10{,}000 & 0.1019  & ${\sim}1$:1 & No \\
\midrule
H3 (H-ch6) & 6 & \textbf{Yes} & 10{,}000 & \textbf{0.00009} & \textbf{70{,}404:1} & \textbf{Yes} \\
\bottomrule
\end{tabular}
\end{table}

The spectral-only experiments converge to a Fiedler band of
$0.0836$--$0.1019$ (Figure~\ref{fig:spectral_attractor})---a 10.4\%
coefficient of variation across a $4\times$ span of channel counts
($n = 3$ to $n = 12$). The
attractor value clusters near the spectral loss target of
$\lambda_2^* = 0.1$ described in Paper~3. This convergence is
\emph{channel-count-independent}: the bridge reaches the same
connectivity equilibrium whether it has 3, 4, 8, or 12~channels.

At the attractor, bridge structure is \emph{connected but
undirected}. Co-planar/cross-planar ratios are ${\sim}1$:1 (where
applicable), and no block-diagonal structure emerges. The spectral
loss distributes connectivity uniformly across all channel pairs;
without contrastive specification, there is no reason for the bridge
to favor any particular pair over another.

H-ch12 ($n = 12$) is the most demanding case. With
$\binom{12}{2} = 66$ off-diagonal pairs sharing a fixed connectivity
budget, convergence is slower but the bridge reaches Fiedler $0.1019$
at step~10{,}000---confirming the attractor holds at $4\times$ the
minimum channel count. The slight upward trend (0.1019 vs.\
$0.0836$--$0.0951$ for smaller~$n$) suggests a weak positive
relationship between channel count and final Fiedler, but the effect
is small relative to the attractor's stability across the full range.

\subsection{Interpretation}

The spectral attractor has a straightforward interpretation. The
spectral loss penalizes deviations of the bridge's Fiedler eigenvalue
from a target~$\lambda_2^*$. Under spectral regularization alone, the
bridge distributes coupling uniformly across all off-diagonal entries:
this is the minimum-energy configuration that achieves the target
$\lambda_2$. Adding more channels does not change this equilibrium
because each new channel adds both off-diagonal entries (coupling
budget to distribute) and the same spectral target (connectivity
to achieve). The ratio of budget to target remains constant.

The attractor is the bridge's \emph{unprogrammed state}---the
resting connectivity it achieves before any topological direction
is specified. It is analogous to a uniform probability distribution:
maximum entropy given the connectivity constraint. The spectral
loss provides the constraint (how much total connectivity); the
contrastive loss provides the structure (where connectivity
concentrates).

\begin{figure}[t]
  \centering
  \includegraphics[width=\textwidth]{figures/fig_spectral_attractor.pdf}
  \caption{Spectral attractor universality across channel counts.
    Fiedler eigenvalue converges to ${\approx}\,0.09$--$0.10$ under
    spectral-only training for $n \in \{3, 4, 8, 12\}$---a 10.4\%
    band across a $4\times$ range of channel counts. The attractor is
    the bridge's natural connectivity equilibrium in the absence of
    contrastive topology specification.}
  \label{fig:spectral_attractor}
\end{figure}

\subsection{The Bifurcation}
\label{sec:bifurcation}

The spectral attractor provides the baseline against which
contrastive topology programming is measured. Adding contrastive
loss breaks the attractor by specifying directional coupling
preferences, producing a sharp bifurcation in spectral connectivity.

The bifurcation is extreme at both confirmed dimensionalities:

\begin{itemize}
  \item \textbf{RD ($n = 6$):} Fiedler drops from $0.0918$
    (spectral-only, H-ch4 at $n = 4$, nearest comparable) to
    $0.00009$ (contrastive active). A factor of $1{,}020\times$.
    Co/cross: ${\sim}1$:1 $\to$ $70{,}404$:1.
  \item \textbf{Tesseract ($n = 8$):} Fiedler drops from $0.0944$
    (spectral-only, H-ch8) to $0.000191$ (contrastive active,
    T-001r2). A factor of ${\sim}470\times$. Co/cross: ${\sim}1$:1
    $\to$ $41{,}564$:1.
  \item \textbf{Block-diagonal classification:} 0\% (spectral-only
    at all channel counts) $\to$ 100\% (contrastive active at
    $n = 6$ and $n = 8$).
\end{itemize}

This bifurcation decomposes the Steersman's effect into two
orthogonal control channels:

\begin{enumerate}
  \item \textbf{Spectral loss $\to$ connectivity.} The spectral
    component drives the bridge toward a target algebraic
    connectivity ($\lambda_2 \approx 0.1$). It makes the bridge
    \emph{connected}---channels interact---but imposes no directional
    preference. The result is the spectral attractor: a connected,
    isotropic bridge.
  \item \textbf{Contrastive loss $\to$ topology.} The contrastive
    component specifies which connections to strengthen (co-axial
    pairs) and which to suppress (cross-axial pairs). It transforms
    isotropic connectivity into directed topology. The result is
    block-diagonal structure whose block count matches the number
    of co-axial pairs.
\end{enumerate}

Neither component alone produces block-diagonal structure. Spectral
loss alone produces the attractor. Contrastive loss alone (without
spectral regularization) was not tested in this series but is a
natural follow-up. The two components together produce a controlled
bifurcation: connectivity plus direction yields topology.

\subsection{Comparison: Bifurcation Across Dimensionalities}

Table~\ref{tab:bifurcation} summarizes the confirmed bifurcation
events. The bifurcation ratio is smaller at $n = 8$ than at $n = 6$,
likely because the larger number of cross-axial entries at $n = 8$
(24~vs.~12) creates more total suppression pressure, producing a
slightly higher residual Fiedler at convergence. Both bifurcations
are of the same order of magnitude ($10^{2\text{--}3}\times$) and
both produce complete block-diagonal structure (100\% classification).

\begin{table}[H]
\centering
\caption{Bifurcation comparison: spectral attractor broken by
  contrastive topology specification at two dimensionalities.}
\label{tab:bifurcation}
\small
\begin{tabular}{lccccc}
\toprule
Geometry & $n$ & Spectral Fiedler & Contrastive Fiedler & Bifurcation & $\rho$ \\
\midrule
RD (3D)       & 6 & 0.0918 & 0.00009  & $1{,}020\times$ & 70{,}404:1 \\
Tesseract (4D) & 8 & 0.0944 & 0.000191 & ${\sim}470\times$  & 41{,}564:1 \\
\bottomrule
\end{tabular}
\end{table}

% ===================================================================
\section{Topology Programmability}
\label{sec:programmability}

Sections~\ref{sec:tesseract} and~\ref{sec:attractor} established that
the Steersman programs polytope-derived topologies and that the
spectral attractor provides the unprogrammed baseline. This section
asks a sharper question: does the Steersman program \emph{any} pair
specification, or does the pair specification require geometric
coherence?

Two negative controls bracket the answer: wrong-labels (random
partition, WL-001) and octahedral contrastive (minimum geometry,
O-001).

\subsection{Wrong-Labels: Random Partition (WL-001)}

The wrong-labels experiment uses $n = 6$ channels with a random pair
partition instead of the RD face-pair geometry. All other parameters
(model, seed, hyperparameters, Steersman configuration) are identical
to the standard contrastive experiments.
If the Steersman is a brute-force topology programmer, the random
partition should produce block-diagonal structure aligned with the
random pairs. If geometric coherence is required, no block structure
should emerge.

\begin{table}[H]
\centering
\caption{Wrong-labels control ($n = 6$, random pair partition) vs.\
  standard contrastive ($n = 6$, RD face pairs). Both experiments use
  identical hyperparameters. The wrong-labels co/cross value reports
  alignment with the \emph{random} pair specification.}
\label{tab:wronglabels}
\small
\begin{tabular}{lccccc}
\toprule
Experiment & Pairs & Steps & Fiedler & $\rho$ & BD \\
\midrule
H-ch6 (RD) & $(0,5), (1,4), (2,3)$ & 10{,}000 & 0.00009 & 70{,}404:1 & Yes \\
WL-001 (random) & random partition & 10{,}000 & $1.26 \times 10^{-5}$ & $8.7 \times 10^{-6}$:1 & \textbf{No} \\
\bottomrule
\end{tabular}
\end{table}

No block-diagonal structure emerged at any point during 10{,}000
steps of wrong-labels training. The co-axial/cross-axial coupling
ratio reached $8.7 \times 10^{-6}$:1---effectively zero. For
comparison, the standard contrastive control (H-ch6) exceeded 200:1
by step~200 and converged to 70{,}404:1.

The Steersman drove Fiedler to $1.26 \times 10^{-5}$---below the
spectral attractor at ${\sim}0.09$---indicating that the contrastive
loss produced \emph{some} structural effect. But the effect was pure
suppression of connectivity, not block-diagonal organization. The
bridge decoupled without organizing: it lost connectivity without
gaining topology.

This is a precise failure mode. The contrastive loss drives down
both co-axial and cross-axial coupling when the pair specification
lacks geometric coherence. Without a coherent spatial decomposition,
there is no gradient signal distinguishing co-axial from cross-axial
entries---the contrastive pressure distributes uniformly and
suppresses all coupling toward zero. The result is a bridge that
is neither connected (Fiedler $\approx 0$) nor structured
(co/cross $\approx 0$)---worse than the unprogrammed attractor.

\subsection{Octahedral Contrastive: Minimum Geometry (O-001)}

The octahedral experiment tests the other boundary: the simplest
possible geometric specification. With $n = 4$ and only 2~co-axial
pairs against 4~cross-axial pairs, the octahedral topology provides
the minimum spatial structure that could plausibly produce
block-diagonal organization.

\begin{table}[H]
\centering
\caption{Octahedral contrastive ($n = 4$, 2 co-axial pairs from
  octahedral vertex opposition) at 10{,}000 steps. Per-bridge mean
  ratio reported; pooled ratio is 369{,}365:1.}
\label{tab:octahedral}
\small
\begin{tabular}{lccccc}
\toprule
Experiment & $n$ & Pairs & Fiedler & $\rho$ & BD \\
\midrule
H-ch4 (spectral-only) & 4 & none & 0.0836 & ${\sim}1$:1 & No \\
O-001 (octahedral) & 4 & $(0,1), (2,3)$ & $1.1 \times 10^{-5}$ & \textbf{473{,}622:1} & \textbf{Yes} \\
\bottomrule
\end{tabular}
\end{table}

O-001 produces the strongest block-diagonal signal in the entire
research programme (Figure~\ref{fig:o001_emergence}): per-bridge mean co-axial/cross-axial coupling
ratio of \textbf{473{,}622:1} (pooled: 369{,}365:1, median:
401{,}851:1, range: 126{,}782--1{,}577{,}518:1 across 88~bridges).
This exceeds both the RD result (70{,}404:1 at $n = 6$) and the
tesseract result (41{,}564:1 at $n = 8$) by an order of magnitude.

\begin{figure}[t]
  \centering
  \includegraphics[width=\textwidth]{figures/fig_o001_emergence.pdf}
  \caption{Octahedral block-diagonal emergence (O-001, $n = 4$).
    Co-axial/cross-axial coupling ratio grows monotonically to
    473{,}622:1---the programme's strongest signal. The minimum
    geometry (two co-axial pairs, four cross-axial pairs) produces
    maximum coupling strength, establishing the inverse-$n$
    relationship.}
  \label{fig:o001_emergence}
\end{figure}

The inverse relationship between channel count and coupling ratio
is established across three data points: $n = 4$ (${\sim}474$K:1)
$>$ $n = 6$ (${\sim}70$K:1) $>$ $n = 8$ (${\sim}42$K:1). Fewer channels means fewer
cross-axial entries to suppress: $n = 4$ has only 4~cross-axial
pairs, while $n = 8$ has 24. The co-axial signal is concentrated
into fewer entries, producing stronger per-pair coupling and a
more extreme ratio. The co-planar coupling magnitude at O-001 also
grows linearly throughout training ($0.086$ at step~500, $1.027$ at
step~10{,}000) with no saturation, indicating the ratio would
increase further with continued training.

\subsection{The Geometric Coherence Requirement}

The wrong-labels and octahedral results, taken together, define the
boundary of topology programmability. On one side: any polytope-derived
spatial specification---RD at $n = 6$, tesseract at $n = 8$, octahedron
at $n = 4$---produces block-diagonal structure with coupling ratios
from 42{,}000:1 to 200{,}000:1. On the other side: a random partition
at the same channel count ($n = 6$) produces no structure at all.

The Steersman is not a brute-force topology enforcer. It is a
\emph{selective geometric compiler}. It compiles valid geometric
programs---pair specifications that encode the co-axial face or vertex
structure of a polytope---and rejects incoherent specifications by
driving the bridge toward total decoupling. The pair specification is
not a suggestion; it is a program that must be well-formed.

This selectivity is a feature, not a limitation. It means the
block-diagonal structure observed under polytope-derived specifications
is genuine geometric encoding, not an artifact of the contrastive loss
formulation. The loss function is identical in WL-001 and H-ch6; only
the pair indices differ. The indices carry the geometric information,
and the Steersman distinguishes coherent from incoherent geometry.

\subsection{Block-Diagonal Hierarchy}

Table~\ref{tab:hierarchy} presents the complete block-diagonal hierarchy
across all experiments, from the strongest positive to the strongest
negative.

\begin{table}[H]
\centering
\caption{Complete topology hierarchy across all experiments, ordered
  by co-axial/cross-axial coupling ratio. Four regimes emerge,
  separated by orders of magnitude.}
\label{tab:hierarchy}
\small
\begin{tabular}{llcccl}
\toprule
Experiment & Topology & $n$ & $\rho$ & Fiedler & Regime \\
\midrule
O-001 & Octahedral & 4 & 473{,}622:1 & $1.1 \times 10^{-5}$ & Geometric BD \\
Seed-43 & RD contrastive & 6 & 73{,}309:1 & ${\sim}10^{-5}$ & Geometric BD \\
H-ch6 & RD contrastive & 6 & 70{,}404:1 & $9 \times 10^{-5}$ & Geometric BD \\
T-001r2 & Tesseract & 8 & 41{,}564:1 & $1.9 \times 10^{-4}$ & Geometric BD \\
\midrule
E-001 & Emanation & 6 & 1.12:1 & 0.084 & Hierarchical \\
\midrule
H-ch3/4/8/12 & Spectral-only & 3--12 & ${\sim}1$:1 & 0.084--0.102 & Spectral attractor \\
\midrule
WL-001 & Wrong-labels & 6 & $8.7 \times 10^{-6}$:1 & $1.3 \times 10^{-5}$ & Collapse \\
R-001 & Resonance & 6 & $8.9 \times 10^{-6}$:1 & $1.2 \times 10^{-5}$ & Collapse \\
\bottomrule
\end{tabular}
\end{table}

The hierarchy separates into four regimes with no overlap
(Figure~\ref{fig:four_regimes}).

\begin{figure}[t]
  \centering
  \includegraphics[width=\textwidth]{figures/fig_four_regimes.pdf}
  \caption{Four dynamical regimes of the bridge matrix under different
    programming conditions, ordered by co-axial/cross-axial coupling
    ratio. Block-diagonal ($\rho > 40{,}000$:1), hierarchical coherence
    ($\rho \approx 1$:1, Fiedler $\approx 0.08$), spectral attractor
    ($\rho \approx 1$:1, Fiedler $\approx 0.09$), and connectivity
    collapse ($\rho \approx 10^{-5}$:1). No intermediate outcomes
    exist between regimes.}
  \label{fig:four_regimes}
\end{figure}

\textbf{Regime 1: Geometric block-diagonal} ($\rho > 40{,}000$:1).
Polytope-derived pair specifications (octahedral, RD, tesseract)
produce extreme co-axial/cross-axial separation with near-zero
Fiedler. The regime spans one order of magnitude internally
(42K--474K) and is separated from the next regime by four orders
of magnitude.

\textbf{Regime 2: Hierarchical coherence} ($\rho \approx 1$:1,
Fiedler $\approx 0.08$). The emanation architecture (E-001) produces
slight directional preference without block-diagonal structure. The
bridge maintains connectivity near the spectral attractor but
develops weak asymmetry through the master-offset hierarchy
(Section~\ref{sec:nongeometric}).

\textbf{Regime 3: Spectral attractor} ($\rho \approx 1$:1,
Fiedler $\approx 0.09$). Spectral-only training (no contrastive
loss) produces connected, isotropic bridges. This is the
unprogrammed state (Section~\ref{sec:attractor}).

\textbf{Regime 4: Connectivity collapse} ($\rho \approx 0$:1,
Fiedler $\approx 0$). Incoherent contrastive specifications
(wrong-labels, resonance) produce total collapse: both
co-axial/cross-axial ratio and Fiedler approach zero. The bridge
disconnects without organizing. This regime is distinct from the
spectral attractor: the spectral attractor maintains stable
connectivity ($\text{Fiedler} \approx 0.09$), while the collapse
regime drives connectivity to zero ($\text{Fiedler} \approx 10^{-5}$).
The contrastive loss, given incoherent pairs, overwhelms the spectral
loss and suppresses all coupling.

The four regimes are cleanly separated. No experiment falls between
regimes. The Steersman's response to a pair specification is
categorical: geometric coherence produces block-diagonal structure;
hierarchical indirection produces weak coherence; no specification
produces the spectral attractor; incoherent specification produces
collapse.

% ===================================================================
\section{Non-Geometric Topologies}
\label{sec:nongeometric}

Section~\ref{sec:programmability} established that the Steersman
requires geometric coherence for block-diagonal structure to emerge.
This section examines the two non-geometric experiments in detail:
circular resonance (R-001, number-theoretic pairs) and the emanation
architecture (E-001, hierarchical bridge structure). Together with
WL-001 (wrong-labels), they define the boundary between successful
and unsuccessful topology programming.

\subsection{Circular Resonance: The Strongest Negative (R-001)}

The resonance topology derives co-axial pair assignments from
number-theoretic relationships: Sophie Germain pairing
($2(11) + 1 = 23 \to$ channels $(0, 2)$), consecutive-prime pairing
($29, 31 \to$ channels $(3, 4)$), and residue pairing
($11 \equiv 17 \equiv 5 \pmod{6} \to$ channels $(1, 5)$). Unlike
the RD pairs, which correspond to opposite faces of a Voronoi cell,
the resonance pairs have algebraic but not spatial motivation.

\begin{table}[H]
\centering
\caption{Resonance control (R-001) vs.\ wrong-labels (WL-001) vs.\
  RD contrastive (H-ch6). All three use $n = 6$ with identical
  hyperparameters; only the pair specification differs.}
\label{tab:resonance}
\small
\begin{tabular}{lcccccc}
\toprule
Experiment & Pair Source & Steps & $\rho$ & Fiedler & Deviation & BD \\
\midrule
H-ch6 (RD) & Spatial & 10{,}000 & 70{,}404:1 & $9 \times 10^{-5}$ & --- & Yes \\
WL-001 (random) & Random & 10{,}000 & $8.7 \times 10^{-6}$ & $1.3 \times 10^{-5}$ & 2.12 & No \\
R-001 (primes) & Algebraic & 10{,}000 & $8.9 \times 10^{-6}$ & $1.2 \times 10^{-5}$ & 2.12 & No \\
\bottomrule
\end{tabular}
\end{table}

R-001 is the strongest negative control in the programme. The
co-axial/cross-axial coupling ratio decayed monotonically from
$0.257$ at step~100 to $6.5 \times 10^{-4}$ at step~1{,}000 and
$8.9 \times 10^{-6}$ at step~10{,}000. Fiedler follows the same
trajectory: $1.7 \times 10^{-4}$ at step~1{,}000, $1.2 \times 10^{-5}$
at step~10{,}000. Deviation from identity grows throughout training
to~2.12, indicating that the bridge matrix departs substantially
from identity without developing any directional structure---an
extreme, unorganized coupling pattern.

Validation loss converges normally ($0.4008$ at step~10{,}000),
confirming that the resonance topology does not impede language
modelling. The adapter learns the task while the bridge fails to
organize.

\subsection{WL-001 and R-001: Functional Equivalence}

A key finding emerges from comparing the two negative controls:
WL-001 (random partition) and R-001 (number-theoretic pairs) produce
\emph{nearly identical trajectories} (Figure~\ref{fig:wl_r001_identity}).
Table~\ref{tab:wl_r_compare} shows matched checkpoints.

\begin{table}[H]
\centering
\caption{Trajectory comparison: WL-001 (random partition) vs.\
  R-001 (number-theoretic pairs). Maximum divergence ${<}\,10\%$
  at all checkpoints.}
\label{tab:wl_r_compare}
\small
\begin{tabular}{ccccc}
\toprule
Step & WL-001 $\rho$ & R-001 $\rho$ & WL-001 Fiedler & R-001 Fiedler \\
\midrule
1{,}000 & $6.5 \times 10^{-4}$ & $7.1 \times 10^{-4}$ & $1.6 \times 10^{-4}$ & $1.7 \times 10^{-4}$ \\
5{,}000 & $5.0 \times 10^{-5}$ & $5.1 \times 10^{-5}$ & $6.0 \times 10^{-5}$ & $6.4 \times 10^{-5}$ \\
10{,}000 & $8.7 \times 10^{-6}$ & $8.9 \times 10^{-6}$ & $1.3 \times 10^{-5}$ & $1.2 \times 10^{-5}$ \\
\bottomrule
\end{tabular}
\end{table}

The maximum divergence between WL-001 and R-001 is less than 10\%
at all checkpoints. From the Steersman's perspective,
number-theoretic structure without spatial embedding is
\emph{functionally equivalent to random noise}. The algebraic
relationships (Sophie Germain, twin primes, residue classes) carry
no information that the contrastive loss can use to produce
directional coupling.

This equivalence has a precise meaning. The contrastive loss
operates on the absolute magnitudes of bridge entries at the
indices specified by the pair partition. For the contrastive gradient
to produce block-diagonal structure, the specified co-axial entries
must be amenable to selective strengthening relative to the
cross-axial entries. When the pair specification reflects the
geometry of a polytope (RD, tesseract, octahedron), the bridge's
gradient landscape has a clear minimum corresponding to the
block-diagonal attractor. When the pair specification is algebraic
(resonance) or random (wrong-labels), no such minimum exists---the
gradient landscape is flat with respect to the co-axial/cross-axial
distinction, and the contrastive pressure distributes uniformly,
driving total collapse.

Both experiments converge to Regime~4 (connectivity collapse):
Fiedler $\approx 10^{-5}$, co/cross $\approx 10^{-5}$, deviation
$\approx 2.1$. The bridge develops extreme but undirected coupling
and then decouples. This is distinct from the spectral attractor
(Regime~3), where the bridge maintains stable connectivity at
Fiedler~$\approx 0.09$. The contrastive loss, given incoherent
pairs, overwhelms the spectral loss's connectivity target and
suppresses all coupling---a destructive interaction between the
two loss components.

\begin{figure}[t]
  \centering
  \includegraphics[width=\textwidth]{figures/fig_wl_r001_identity.pdf}
  \caption{Functional equivalence of wrong-labels (WL-001, random
    partition) and resonance (R-001, number-theoretic pairs). Both
    collapse trajectories are indistinguishable ($<10\%$ divergence
    across 10{,}000 steps), demonstrating that algebraic structure
    without spatial embedding carries no information the contrastive
    loss can use.}
  \label{fig:wl_r001_identity}
\end{figure}

\subsection{Emanation Architecture: The Third Regime (E-001)}
\label{sec:emanation_results}

The emanation architecture introduces a qualitatively different
outcome. Using a master bridge matrix with per-layer offsets scaled
by $1/\sqrt{l}$ (Eq.~\ref{eq:emanation}), the same RD-derived
co-axial pairs are specified, but the contrastive loss operates
through the master matrix rather than per-layer bridges directly.

\begin{table}[H]
\centering
\caption{Emanation results (E-001) at 10{,}000 steps vs.\ standard
  RD contrastive (H-ch6) and spectral attractor (H-ch3).}
\label{tab:emanation}
\small
\begin{tabular}{lcccccl}
\toprule
Experiment & Architecture & $\rho$ & Fiedler & Deviation & Val Loss & Regime \\
\midrule
H-ch6 (RD) & Independent & 70{,}404:1 & $9 \times 10^{-5}$ & --- & 0.4015 & BD \\
E-001 (emanation) & Hierarchical & 1.12:1 & 0.084 & 0.35 & 0.4009 & Hierarchical \\
H-ch3 (spectral) & Independent & N/A & 0.095 & --- & 0.4020 & Attractor \\
\bottomrule
\end{tabular}
\end{table}

E-001 occupies a regime distinct from both block-diagonal emergence
and the spectral attractor. Its final metrics at 10{,}000 steps:

\begin{itemize}
  \item \textbf{Fiedler: 0.084.} Near but below the spectral attractor
    ($0.09$--$0.10$). The bridge maintains substantial connectivity.
  \item \textbf{Co/cross: 1.12:1.} Slight directional preference---the
    bridge favors the RD co-axial pairs by 12\% on average. This is
    above unity but far below the 70{,}000:1 achieved by the standard
    architecture with the same pair specification.
  \item \textbf{Deviation: 0.35.} The per-layer offsets have grown
    substantially from their initialized values, indicating that
    individual layers differentiate from the master bridge. This
    deviation is roughly one-sixth the deviation observed in the
    collapse regime (2.12 for WL-001/R-001).
  \item \textbf{Val loss: 0.4009.} Marginally better than both the
    standard contrastive (0.4015) and spectral-only (0.4020)
    experiments, though the difference is within noise.
\end{itemize}

\paragraph{Interpretation.}
The emanation architecture creates a tension between global coherence
and local adaptation. The master bridge receives the contrastive
gradient directly, developing slight co-axial preference. But the
per-layer offsets dilute this preference: each layer adapts its bridge
to local gradient signals, and these local adaptations do not
consistently reinforce the global topology. The master bridge carries
a weak geometric signal (12\% co-axial preference); the offsets
fragment it.

The result is a bridge that is \emph{coherent but not structured}:
it maintains connectivity (Fiedler near the attractor), develops weak
directional preference (co/cross slightly above unity), and allows
per-layer differentiation (growing deviation)---without achieving the
dramatic block-diagonal separation that requires per-bridge
contrastive specificity.

This hierarchical coherence regime suggests that topology
programming requires direct contrastive feedback to individual
bridge matrices. Indirection through a shared master bridge preserves
connectivity but loses the topological specificity needed for
block-diagonal emergence. The emanation architecture may have value
for other purposes---regularization, transfer, compositional
adapters---but it does not serve as a topology programmer.

\subsection{Topology Taxonomy: Four Regimes}

The complete experimental programme produces four distinct dynamical
regimes, cleanly separated with no intermediate or overlapping
outcomes.

\begin{enumerate}
  \item \textbf{Block-diagonal} (geometric contrastive).
    Polytope-derived pair specifications with per-bridge contrastive
    loss. Co/cross: $42{,}000$--$474{,}000$:1. Fiedler: $10^{-5}$--$10^{-4}$.
    Experiments: O-001, H-ch6, Seed-43, T-001r2.

  \item \textbf{Spectral attractor} (no contrastive).
    Spectral regularization alone, no pair specification.
    Co/cross: ${\sim}1$:1. Fiedler: $0.092$--$0.102$.
    Experiments: H-ch3, H-ch4, H-ch8, H-ch12.

  \item \textbf{Hierarchical coherence} (emanation architecture).
    Contrastive loss on shared master bridge with per-layer offsets.
    Co/cross: $1.12$:1. Fiedler: $0.084$.
    Experiment: E-001.

  \item \textbf{Connectivity collapse} (incoherent contrastive).
    Non-spatial pair specifications (random or algebraic).
    Co/cross: ${\sim}10^{-5}$:1. Fiedler: ${\sim}10^{-5}$.
    Deviation: ${\sim}2.1$.
    Experiments: WL-001, R-001.
\end{enumerate}

Two features of this taxonomy are notable. First, the collapse regime
(Regime~4) is distinct from the spectral attractor (Regime~3). The
spectral attractor is a near-initialization resting state at Fiedler
$\approx 0.09$---inside the BM-000 identity-plus-noise null, not a
distinct structured equilibrium; the collapse regime represents a
destructive interaction where the contrastive loss overwhelms the
spectral loss and drives connectivity to zero. The spectral attractor
is the bridge's resting state; the collapse regime is the bridge's
failure mode under incoherent programming.

Second, Regimes~2 and~3 are structurally similar (both near
Fiedler~$\approx 0.09$, both with co/cross near unity) but
mechanistically different. The spectral attractor sits at
near-initialization connectivity under spectral loss alone.
Hierarchical coherence reaches a similar Fiedler through a different
pathway---master bridge with contrastive loss---and develops weak
directionality that the spectral attractor does not. The regimes are adjacent in
metric space but distinct in mechanism.

% ===================================================================
\section{Discussion}
\label{sec:discussion}

\subsection{The Bridge as Programmable Substrate}

Paper~3 framed the result as \emph{discovery}: the Steersman
discovers rhombic dodecahedral geometry in the bridge matrix. The
present paper reframes this as \emph{programming}: the Steersman
programs whatever geometry it is told, provided the pair
specification corresponds to a coherent spatial structure.

The evidence supports the reframing. The tesseract result
(Section~\ref{sec:tesseract}) demonstrates that the Steersman
programs 4D geometry with the same fidelity as 3D: four independent
$2 \times 2$ blocks, clean eigenvalue split, coupling ratio
exceeding 41{,}000:1. The octahedral result (O-001) demonstrates
programming at minimum dimensionality, achieving the programme's
strongest block-diagonal signal (473{,}622:1 per-bridge mean). The
spectral attractor (Section~\ref{sec:attractor}) shows that the
bridge has a natural unprogrammed state---connected, isotropic,
Fiedler~$\approx 0.09$---and that contrastive loss adds direction
to this connectivity.

The four-regime taxonomy (Section~\ref{sec:programmability}) sharpens
the framing further. The bridge is a programmable substrate, but the
programming language is \emph{spatial}. The Steersman compiles
polytope-derived pair specifications into block-diagonal structure.
It rejects non-spatial specifications---random partitions and
number-theoretic relationships---by driving the bridge into
connectivity collapse. The selectivity is a feature: it confirms
that successful topology programming encodes genuine geometric
information, not an artifact of the loss formulation.

\subsection{The Selective Geometric Compiler}

The central thesis upgrade from Paper~3 is precise. The Steersman
is not a universal topology enforcer---it does not program arbitrary
pair specifications into block structure. It is a \emph{selective
geometric compiler}: it accepts pair specifications derived from
the co-axial face or vertex structure of polytopes and rejects
all others.

This selectivity is demonstrated by the four-way contrast at $n = 6$:

\begin{itemize}
  \item \textbf{RD pairs (H-ch6):} 70{,}404:1. Spatial
    specification. Block-diagonal.
  \item \textbf{Emanation (E-001):} 1.12:1. Same spatial pairs,
    indirect application through master bridge. Weak coherence but
    no block-diagonal.
  \item \textbf{Wrong-labels (WL-001):} $8.7 \times 10^{-6}$:1.
    Random pairs. Connectivity collapse.
  \item \textbf{Resonance (R-001):} $8.9 \times 10^{-6}$:1.
    Number-theoretic pairs. Connectivity collapse.
\end{itemize}

The selectivity has two dimensions. The first is geometric: the pair
specification must encode a coherent spatial decomposition (Regimes~1
vs.~4). The second is architectural: the contrastive signal must reach
individual bridge matrices directly, not through hierarchical
indirection (Regimes~1 vs.~2).

\subsection{Number-Theoretic Structure Is Insufficient}

The functional equivalence of WL-001 and R-001 (maximum divergence
${<}\,10\%$ across 10{,}000 steps) has a clear implication:
\emph{mathematical structure without spatial embedding carries no
information that the contrastive loss can use}. The Sophie Germain
relationship ($2(11) + 1 = 23$), the twin prime pair $(29, 31)$, and
the residue class partition are genuine algebraic structures. They
produce the same outcome as random noise.

This result constrains the class of viable topology specifications.
The pair specification must encode a property of the bridge's
gradient landscape, not merely a mathematical relationship among
channel indices. Polytope co-axial structure encodes such a
property: it identifies axis-aligned subspaces of the bridge matrix
that can be independently optimized. Number-theoretic relationships
do not identify such subspaces.

\subsection{Spectral-Contrastive Decomposition}

The decomposition of the Steersman into spectral (connectivity) and
contrastive (topology) components has a precise analogue in graph
theory. The Fiedler eigenvalue measures algebraic
connectivity~\cite{fiedler1973}---how well-connected the graph is.
The contrastive loss specifies edge weights---which connections are
strengthened. Spectral regularization without contrastive
specification produces an Erd\H{o}s--R\'enyi-like coupling
pattern (all edges approximately equal); adding contrastive
specification produces a stochastic block model~\cite{holland1983}
with known community structure.

The four regimes map onto this decomposition:

\begin{itemize}
  \item \textbf{Spectral attractor:} Spectral loss alone. Connectivity
    without topology. The bridge graph is an approximately regular graph.
  \item \textbf{Block-diagonal:} Spectral + coherent contrastive. The
    bridge graph is a disjoint union of cliques (the blocks).
  \item \textbf{Hierarchical coherence:} Spectral + indirectly applied
    contrastive. The bridge graph is weakly modular---slight clustering
    without full community separation.
  \item \textbf{Collapse:} Spectral + incoherent contrastive.
    Destructive interference between the two losses drives the graph
    toward the empty graph.
\end{itemize}

This decomposition suggests that the two loss components are
independently tunable but interact non-linearly. When they are
aligned (coherent contrastive + spectral), they produce strong
structure. When they conflict (incoherent contrastive + spectral),
the contrastive loss dominates and suppresses connectivity below the
spectral target. Understanding this interaction is a prerequisite
for designing more sophisticated topology specifications.

\subsection{The Inverse-$n$ Relationship}

The coupling ratio hierarchy---O-001 ($n = 4$, ${\sim}474$K:1) $>$
H-ch6 ($n = 6$, ${\sim}70$K:1) $>$ T-001r2 ($n = 8$,
${\sim}42$K:1)---follows a clear inverse relationship with channel
count (Figure~\ref{fig:inverse_n}). The mechanism is straightforward:
at higher~$n$, the bridge
has more cross-axial entries to suppress, diluting the contrastive
signal per entry. At $n = 4$, only 4~cross-axial pairs compete
with 2~co-axial pairs; at $n = 8$, 24~cross-axial pairs compete
with 4~co-axial pairs. The co-axial/cross-axial pair count ratio
($|\mathcal{C}|/|\mathcal{X}|$) is $2/4 = 0.50$ for octahedral,
$3/12 = 0.25$ for RD, and $4/24 = 0.167$ for tesseract.

This suggests a practical guideline: \emph{use the minimum channel
count that admits the desired topology}. Higher channel counts
provide more topological expressiveness (more independent axes) but
weaker per-axis signal. The practitioner's choice of~$n$ trades off
topological resolution against coupling strength.

\begin{figure}[t]
  \centering
  \includegraphics[width=0.7\textwidth]{figures/fig_inverse_n.pdf}
  \caption{Inverse-$n$ relationship between channel count and
    co-axial/cross-axial coupling ratio. Octahedral ($n = 4$):
    473{,}622:1. RD ($n = 6$): 70{,}404:1. Tesseract ($n = 8$):
    41{,}564:1. Fewer channels concentrate the co-axial signal into
    fewer entries, producing stronger per-pair coupling.}
  \label{fig:inverse_n}
\end{figure}

\subsection{Implications for Adapter Design}

If the bridge is a programmable topology substrate, adapter design
acquires a new degree of freedom. Current practice treats the rank
dimension as homogeneous. Multi-channel extensions introduce
structure but leave it to the optimizer to discover whatever
coupling pattern the task prefers. The Steersman enables
\emph{prescribed} coupling: the practitioner specifies which
channels should interact, and the feedback mechanism enforces the
specification during training.

\begin{itemize}
  \item \textbf{Topology as hyperparameter.} The pair specification
    is now a hyperparameter alongside rank, learning rate, and target
    modules. Different polytope geometries may suit different task
    families.
  \item \textbf{Compositional adapters.} Adapters with known
    block-diagonal structure compose more predictably than dense
    adapters: blocks can be independently swapped, scaled, or merged.
    The $2 \times 2$ block is the natural unit of composition.
  \item \textbf{Multi-modal bridges.} When adapting models across
    modalities (text, image, audio), the pair specification could
    encode known cross-modal relationships in the bridge structure.
  \item \textbf{Diagnostic topology.} Train with $n = 6$ or $n = 8$
    and cybernetic feedback to discover the bridge's preferred
    topology, then deploy with lower~$n$ for parameter efficiency.
    Paper~3 showed that $n = 3$ matches $n = 6$ task performance
    with $4\times$ fewer bridge parameters.
\end{itemize}

\subsection{Zero-Cost Bridge: The P0 Proof}

A practical question precedes all topology programming results: does
the learnable bridge degrade task performance? If the $n^2$ bridge
parameters and Steersman feedback impose a cost on the underlying
fine-tuning objective, topology programming would face a performance
tax before its benefits could be assessed.

We test this directly. On Qwen~2.5-1.5B-Instruct---a different model
family and scale from the TinyLlama experiments above---we train two
configurations with identical hyperparameters (rank~24, seed~42,
alpaca-cleaned, cosine LR schedule): standard LoRA and TeLoRA with a
learnable bridge ($n = 6$, identity initialization, no contrastive
loss).

\begin{table}[H]
\centering
\caption{P0 head-to-head: standard LoRA vs.\ TeLoRA with learnable
  bridge (no contrastive loss). Qwen~2.5-1.5B-Instruct, $r = 24$,
  seed~42, 3{,}235~steps.}
\label{tab:p0}
\small
\begin{tabular}{lcccc}
\toprule
Configuration & Final Loss & Trainable Params & Bridge Params & Time \\
\midrule
Standard LoRA ($r = 24$) & 0.1763 & 3{,}268{,}608 & --- & 42.8 min \\
TeLoRA ($r = 24$, $n = 6$) & 0.1762 & 3{,}270{,}624 & 2{,}016 (0.06\%) & 58.3 min \\
\bottomrule
\end{tabular}
\end{table}

Loss is identical to four significant figures (0.01\% difference,
within noise). The bridge adds 2{,}016~parameters---36~per bridge
matrix across 56~target projections---representing 0.06\% of the
trainable parameter budget. The 36.2\% training time overhead
reflects the Steersman's per-step Fiedler eigenvalue computation
and control law evaluation.

Without contrastive loss, the bridges settle near the spectral
attractor: mean Fiedler~$= 0.038$ (somewhat below the
TinyLlama attractor at~$0.09$, consistent with model-dependent
variation), co-axial/cross-axial ratio~$\approx 1$:1 (isotropic,
no directional preference), and mean deviation from identity~$= 0.07$
(Figure~\ref{fig:p0_proof}). The bridge evolves under spectral
pressure alone but maintains the unprogrammed baseline---exactly
the attractor behavior documented in Section~\ref{sec:attractor}.

\begin{figure}[t]
  \centering
  \includegraphics[width=\textwidth]{figures/fig_p0_proof.pdf}
  \caption{P0 head-to-head comparison on Qwen~2.5-1.5B-Instruct.
    Top left: loss curves overlap exactly.
    Top right: training time comparison (36.2\% overhead from
    Steersman computation). Middle: bridge matrix evolution at
    steps~500, 2{,}000, and~3{,}235---identity-like structure with
    slight evolution but no block-diagonal emergence (no contrastive
    loss applied). Bottom: Fiedler connectivity, co/cross ratio, and
    deviation from identity over training.}
  \label{fig:p0_proof}
\end{figure}

The P0 result establishes that the learnable bridge is a zero-cost
container. It adds negligible parameters, preserves task performance,
and defaults to the spectral attractor when no topology is specified.
Topology programming via the Steersman adds structure to this
container at no additional parameter cost beyond the 36~entries per
bridge.

\subsection{Limitations}

Several limitations qualify the present results.

\textbf{Single task family.} All experiments use instruction-following
fine-tuning on alpaca-cleaned. Whether the four-regime taxonomy,
spectral attractor, and topology programming transfer to other task
types (classification, summarization, code generation) remains
untested.

\textbf{Pair-based specifications only.} The current pair
specification interface partitions channels into pairs. More complex
topologies---cliques of size~${>}\,2$, hierarchical nesting, cyclic
structures---would require extending the contrastive loss
formulation.

\textbf{No formal optimality proof.} We demonstrate that the Steersman
programs a given topology, but we do not prove that any particular
topology is optimal for any downstream task. The topology is a
hyperparameter; its relationship to task performance remains to be
characterized.

\textbf{No theoretical account of the attractor.} The spectral
attractor ($\text{Fiedler} \approx 0.09$) is empirically consistent
across channel counts, but we provide no theoretical derivation
of the equilibrium value. A theoretical account, possibly through
random matrix theory or mean-field analysis of the spectral loss
dynamics, is future work.

\textbf{No theoretical account of geometric selectivity.} We observe
that only polytope-derived pair specifications produce block-diagonal
structure, but we do not provide a formal characterization of which
pair specifications constitute ``valid geometric programs.''
Understanding why the contrastive loss landscape has minima at
polytope-aligned configurations and not at algebraically structured
but spatially unembedded configurations is an open question.

\textbf{Single seed for most experiments.} All experiments use random
seed~42 except Seed-43 and Seed-44 (which use seeds~43 and~44
respectively). While the binary regime classification is robust given
the magnitude of the effects, specific numerical quantities---coupling
ratios, Fiedler values, deviation---are point estimates.

\textbf{Emanation architecture incompletely explored.} E-001 tests one
hierarchical design (master bridge + $1/\sqrt{l}$-scaled offsets). The
space of hierarchical bridge architectures is large; other designs may
produce different outcomes. The third regime (hierarchical coherence)
is characterized by a single experiment.

\textbf{Scale.} All topology experiments use TinyLlama-1.1B. Paper~3
established scale consistency for RD contrastive at 1.1B and 7B;
the generalization of 4D programming, geometric selectivity, and the
four-regime taxonomy to larger models is untested.



% ===================================================================
\section{Conclusion}
\label{sec:conclusion}

The bridge matrix in a multi-channel LoRA adapter is a programmable
topology substrate. The Steersman---a cybernetic feedback mechanism
comprising three closed-loop control laws---compiles polytope-derived
pair specifications into block-diagonal structure with coupling ratios
exceeding 40{,}000:1, preserving the property across 3D (rhombic
dodecahedron, octahedron) and 4D (tesseract) geometries.

The mechanism is confirmed at three dimensionalities and four channel
counts:

\begin{itemize}
  \item \textbf{Octahedron} ($n = 4$, 3D): two $2 \times 2$ blocks,
    per-bridge mean coupling ratio \textbf{473{,}622:1}. The
    programme's strongest signal.
  \item \textbf{Rhombic dodecahedron} ($n = 6$, 3D): three
    $2 \times 2$ blocks, coupling ratio \textbf{70{,}404:1}. Paper~3's
    result, reproduced across seeds (Seed-43: 73{,}309:1).
  \item \textbf{Tesseract} ($n = 8$, 4D): four $2 \times 2$ blocks,
    coupling ratio \textbf{41{,}564:1}. Clean $4{+}4$ eigenvalue
    split. Reproducible ($r = 1.0000$).
\end{itemize}

Three findings sharpen the understanding of this mechanism.
\emph{First}, the spectral attractor at Fiedler ${\approx}\,0.09$
is universal across channel counts ($n \in \{3, 4, 8, 12\}$, 10.4\%
band), providing the natural connectivity equilibrium from which
contrastive programming departs. \emph{Second}, the Steersman is a
selective geometric compiler: it accepts spatial pair specifications
and rejects algebraic ones, as demonstrated by the functional
equivalence of wrong-label and resonance controls ($<10\%$ trajectory
divergence over 10{,}000 steps). \emph{Third}, the inverse-$n$
relationship---474K:1 ($n = 4$) $>$ 70K:1 ($n = 6$) $>$ 42K:1
($n = 8$)---establishes a practical design principle: minimum
geometry produces maximum signal.

The four-regime taxonomy---block-diagonal, spectral attractor,
hierarchical coherence, and connectivity collapse---provides a
complete map of the bridge's dynamical behavior under different
programming conditions. These regimes are cleanly separated with no
intermediate outcomes, suggesting that the spectral-contrastive
decomposition is a fundamental feature of the loss landscape.

Topology is now a hyperparameter. The choice of polytope geometry
joins rank, learning rate, and target modules as a design decision
for adapter construction. The pair specification is the program; the
Steersman is the compiler; the bridge is the substrate. The compiler
accepts valid geometric programs and rejects incoherent ones---and
this selectivity confirms that successful topology programming
encodes genuine geometric information, not an artifact of the loss
formulation.


% ===================================================================
\appendix
\section{Experimental Details}
\label{app:details}

\subsection{Model and Training Configuration}

All experiments use TinyLlama-1.1B-Chat-v1.0 with the following
hyperparameters held constant across all experiments:

\begin{table}[H]
\centering
\caption{Shared training configuration across all experiments.}
\label{tab:config}
\small
\begin{tabular}{ll}
\toprule
Parameter & Value \\
\midrule
Base model & TinyLlama-1.1B-Chat-v1.0 \\
Dataset & alpaca-cleaned~\cite{alpaca} \\
Training steps & 10{,}000 \\
Learning rate & $2 \times 10^{-4}$ \\
LR scheduler & Cosine \\
Batch size & 2 \\
Gradient accumulation & 8 (effective batch 16) \\
LoRA rank ($r$) & 24 \\
LoRA alpha & 32 \\
LoRA target modules & q\_proj, k\_proj, v\_proj, o\_proj \\
Initialization & Identity \\
Max sequence length & 512 \\
Precision & bfloat16 \\
Random seed & 42 (except Seed-43: 43, Seed-44: 44) \\
\bottomrule
\end{tabular}
\end{table}

The channel size $s = r / n$ varies with channel count: $s = 8$
($n = 3$), $s = 6$ ($n = 4$), $s = 4$ ($n = 6$), $s = 3$ ($n = 8$),
$s = 2$ ($n = 12$). All experiments run to 10{,}000 steps with
checkpoints saved every 100 steps for the first 1{,}000 steps and
every 500 steps thereafter.

\subsection{Steersman Configuration}

The Steersman applies two loss components to the bridge matrix
$\mathbf{M} \in \mathbb{R}^{n \times n}$:

\paragraph{Spectral loss.}
Regularizes the Fiedler eigenvalue $\lambda_2$ of the weighted graph
Laplacian $\mathbf{L} = \mathbf{D} - |\mathbf{M}|$, where
$\mathbf{D}$ is the degree matrix of $|\mathbf{M}|$. This loss
encourages algebraic connectivity without specifying direction.

\paragraph{Contrastive loss.}
Given a set of co-axial pairs $\mathcal{C}$ and the complementary set
of cross-axial pairs $\mathcal{X}$, the contrastive loss rewards
$|\mathbf{M}_{ij}|$ for $(i,j) \in \mathcal{C}$ and penalizes
$|\mathbf{M}_{ij}|$ for $(i,j) \in \mathcal{X}$.

\paragraph{Control laws.}
The Steersman monitors three diagnostics: (1)~connectivity trending
(Fiedler eigenvalue trajectory), (2)~directionality stagnation
(co-axial/cross-axial coupling ratio plateau detection), and
(3)~stability monitoring (loss component variance). These three
scalar diagnostics modulate the relative weighting of spectral and
contrastive losses. The control laws are identical across all
experiments; only the pair specification $\mathcal{C}$ changes.

\subsection{Experiment Index}

\begin{table}[H]
\centering
\caption{Complete experiment index. All experiments use identical
  Steersman configuration and training hyperparameters; only $n$,
  pair specification, and architecture differ.}
\label{tab:experiments}
\small
\begin{tabular}{llcclll}
\toprule
ID & Topology & $n$ & $|\mathcal{C}|$ & Pairs & Architecture & Section \\
\midrule
O-001 & Octahedral & 4 & 2 & $(0,1),(2,3)$ & Independent & \ref{sec:programmability} \\
H-ch6 & RD contrastive & 6 & 3 & $(0,5),(1,4),(2,3)$ & Independent & \ref{sec:programmability} \\
Seed-43 & RD contrastive & 6 & 3 & $(0,5),(1,4),(2,3)$ & Independent & \ref{sec:programmability} \\
Seed-44 & RD contrastive & 6 & 3 & $(0,5),(1,4),(2,3)$ & Independent & \ref{sec:programmability} \\
T-001r2 & Tesseract & 8 & 4 & $(0,1),(2,3),(4,5),(6,7)$ & Independent & \ref{sec:tesseract} \\
\midrule
H-ch3 & Spectral-only & 3 & --- & --- & Independent & \ref{sec:attractor} \\
H-ch4 & Spectral-only & 4 & --- & --- & Independent & \ref{sec:attractor} \\
H-ch8 & Spectral-only & 8 & --- & --- & Independent & \ref{sec:attractor} \\
H-ch12 & Spectral-only & 12 & --- & --- & Independent & \ref{sec:attractor} \\
\midrule
E-001 & Emanation & 6 & 3 & $(0,5),(1,4),(2,3)$ & Hierarchical & \ref{sec:nongeometric} \\
\midrule
WL-001 & Wrong-labels & 6 & 3 & Random partition & Independent & \ref{sec:programmability} \\
R-001 & Resonance & 6 & 3 & $(0,2),(3,4),(1,5)$ & Independent & \ref{sec:nongeometric} \\
\midrule
P0-LoRA & Standard LoRA & --- & --- & --- & Baseline & \ref{sec:discussion} \\
P0-TeLoRA & TeLoRA (no contrastive) & 6 & --- & --- & Independent & \ref{sec:discussion} \\
\bottomrule
\end{tabular}
\end{table}

\subsection{Metrics}

\paragraph{Co-axial/cross-axial coupling ratio ($\rho$).}
For each bridge matrix $\mathbf{M}$, compute the mean absolute value
of co-axial entries (those at positions $(i,j) \in \mathcal{C}$) and
cross-axial entries (those at positions $(i,j) \in \mathcal{X}$). The
ratio $\rho = \bar{|M_\mathcal{C}|} \,/\, \bar{|M_\mathcal{X}|}$
measures directional preference. Values $\gg 1$ indicate
block-diagonal structure; values $\approx 1$ indicate isotropy;
values $\ll 1$ indicate suppression of the specified co-axial
direction.

\paragraph{Fiedler eigenvalue ($\lambda_2$).}
The second-smallest eigenvalue of the weighted graph Laplacian
$\mathbf{L} = \mathbf{D} - |\mathbf{M}|$. Measures algebraic
connectivity of the bridge. Values near zero indicate
near-disconnection (block-diagonal or collapse); values near $0.1$
indicate stable connectivity (spectral attractor).

\paragraph{Deviation from identity.}
Mean Frobenius norm of the difference between each bridge matrix and
the identity matrix, normalized by dimension. Measures how far the
bridge has departed from initialization. Reported primarily for the
collapse and emanation regimes.

\paragraph{Eigenvalue split.}
For $n$-channel bridges with $k$ co-axial pairs, successful topology
programming produces a $k + (n-k)$ eigenvalue split: $k$~eigenvalues
elevated (within-block connectivity) and $n-k$~eigenvalues suppressed
to $\mathcal{O}(10^{-4})$ (between-block disconnection).

\subsection{Hardware}

Experiments were conducted on two machines:

\begin{itemize}
  \item \textbf{Local:} NVIDIA RTX 6000 Ada Generation (48\,GB VRAM).
    Used for T-001r2 (tesseract contrastive).
  \item \textbf{Hermes:} NVIDIA RTX 4090 (16\,GB VRAM). Used for
    spectral-only ablation (H-ch3, H-ch4, H-ch8, H-ch12),
    wrong-labels (WL-001), resonance (R-001), emanation (E-001),
    octahedral (O-001), and seed replication (Seed-43, Seed-44).
\end{itemize}

The P0 head-to-head experiment (Section~\ref{sec:discussion}) uses
Qwen~2.5-1.5B-Instruct on the local RTX~6000~Ada to validate
zero-cost bridge overhead on a different model family.


% ===================================================================
\section{Supplementary Figures}
\label{app:supplementary}

\begin{figure}[H]
  \centering
  \includegraphics[width=\textwidth]{figures/fig_24cell_emergence.pdf}
  \caption{24-cell topology programming ($n = 24$). The 24-cell is a
    self-dual regular 4-polytope with 24~vertices forming 12~co-axial
    pairs. Co-axial/cross-axial coupling ratio reaches 35{,}808:1,
    extending the inverse-$n$ relationship beyond the tesseract
    ($n = 8$) to higher channel counts. The block-diagonal attractor
    generalizes to 12~independent $2 \times 2$ blocks.}
  \label{fig:24cell_emergence}
\end{figure}

\begin{figure}[H]
  \centering
  \includegraphics[width=\textwidth]{figures/fig_fi002_init_independence.pdf}
  \caption{Initialization independence (FI-002). Trained bridge sign
    matrices converge to 100\% identical patterns from initializations
    with Hamming distance 9/12. Final bridge Frobenius distance
    ${\sim}10^{-4}$, confirming that the block-diagonal attractor is
    an initialization-independent fixed point of the Steersman
    dynamics.}
  \label{fig:fi002_init_independence}
\end{figure}

\begin{figure}[H]
  \centering
  \includegraphics[width=\textwidth]{figures/fig_fi003_homeostasis.pdf}
  \caption{Persistence test (FI-003). When the Steersman is disabled
    after topology has emerged, block-diagonal structure dissolves:
    co-axial/cross-axial ratio decays from 12{,}586:1 to
    ${\sim}10$:1 within 1{,}200~steps. The block-diagonal attractor
    is a maintained fixed point requiring continuous cybernetic
    feedback, not a self-sustaining equilibrium.}
  \label{fig:fi003_homeostasis}
\end{figure}

\begin{figure}[H]
  \centering
  \includegraphics[width=\textwidth]{figures/fig_fi004_annealing.pdf}
  \caption{Contrastive weight annealing (FI-004). Coupling ratio as
    a function of fixed contrastive weight $c_w$. Peak performance
    at $c_w = 0.017$ (18{,}671:1). The curve characterizes the
    sensitivity of block-diagonal emergence to the contrastive loss
    magnitude, informing the adaptive Steersman's operating range.}
  \label{fig:fi004_annealing}
\end{figure}


% ===================================================================
\section*{Acknowledgements}

The rhombic library is open source under MPL-2.0. We thank the
reviewers of Papers~1--3 for feedback that shaped this experimental
programme.


% ===================================================================
% --- references ---
\bibliographystyle{plainnat}
\bibliography{paper4}

\end{document}
